% French PDF pp. 25--30, completing the last sentence on p. 31.
% The initial starred note and editorial note (0) are outside the sequence.
\setcounter{footnote}{21}
\renewcommand{\thefootnote}{(\arabic{footnote})}

\noindent
(C 0) Let \(\{F_i\to F\}\) be a family of morphisms of \(\widehat{\mathcal C}\).
If for every representable base change \(S\to F\), the family
\(\{F_i\times_F S\to S\}\) is covering, then the original family is covering too.

(C 1) For every covering family \(\{F_i\to F\}\) and every morphism \(G\to F\),
the family \(\{F_i\times_F G\to G\}\) is covering (\og stability under base change\fg).

(C 2) If \(\{F_i\to F\}\) is a covering family and if, for each \(i\),
\(\{F_{ij}\to F_i\}\) is a covering family, then the composite family
\(\{F_{ij}\to F\}\) is covering (\og stability under composition\fg).

(C 3) If \(\{G_j\to F\}\) is a covering family, and if \(\{F_i\to F\}\)
is a family of morphisms with target \(F\) such that for each \(j\) there exists an
\(i\) such that \(G_j\to F\) factors through \(F_i\to F\), then
\(\{F_i\to F\}\) is covering (\og saturation\fg).

(C 4) Every family consisting of a single isomorphism is covering.

Note that (C 2) and (C 3) also imply:

(C 5) If \(\{F_i\to F\}\) is a family of morphisms with target \(F\) such that
there exists a covering family \(\{G_j\to F\}\) such that for every \(j\)
the family \(\{F_i\times_F G_j\to G_j\}\) is covering, then the family
\(\{F_i\to F\}\) is covering (\og a locally covering family is covering\fg).

\numpara{4.2.4}\label{IV.4.2.4}
Conversely, let \(\mathcal C\) be a category possessing fibered products and,
for each \(S\in\Ob\mathcal C\), let us give ourselves a set of families of
morphisms of \(\mathcal C\) with target \(S\), called covering families,
this datum satisfying axioms (C 1) to (C 4) (and hence also (C 5), which is
a consequence of them). For every \(S\in\Ob\mathcal C\), let \(J(S)\) be
the set of sieves of \(S\) possessing a covering basis (or, what amounts to
the same by (C 3), all of whose bases are covering). Then \(S\mapsto J(S)\)
defines a topology on \(\mathcal C\). The two preceding constructions are
inverse to one another.

In fact, in applications it is not very practical to consider all covering
families, for one sometimes possesses fairly simple descriptions of a
\og sufficient\fg{} number of these families. This leads to the following definitions.

\begin{definition}{4.2.5.0}\label{IV.4.2.5.0}
\nde{This definition (placed in the original after 4.2.5) has been put here,
since it will be used in 6.2.1 in a slightly more general setting than that of 4.2.5.}
Let \(\mathcal C\) be a category. Suppose that, for each \(S\in\Ob\mathcal C\),
a set \(P(S)\) of families of morphisms of \(\mathcal C\) with target \(S\)
is given. One calls the \emph{topology generated by \(P\)} the least fine
topology for which the given families are covering.
\end{definition}

\begin{definition}{4.2.5}\label{IV.4.2.5}
Let \(\mathcal C\) be a category. One calls a \emph{pretopology} on \(\mathcal C\)
the datum, for each \(S\in\Ob\mathcal C\), of a set \(R(S)\) of families
of morphisms \(\{S_i\to S\}\) with target \(S\), called covering for the
pretopology in question, satisfying the following axioms:

(P 1) For every family \(\{S_i\to S\}\in R(S)\) and every morphism \(T\to S\),
the fibered products \(S_i\times_S T\) exist and
\(\{S_i\times_S T\to T\}\in R(T)\).

(P 2) If \(\{S_i\to S\}\in R(S)\) and if for each \(i\),
\(\{T_{ij}\to S_i\}\in R(S_i)\), then the composite family
\(\{T_{ij}\to S\}\) belongs to \(R(S)\).

(P 3) Every family consisting of a single isomorphism is covering.
\end{definition}

\begin{proposition}{4.2.6}\label{IV.4.2.6}
For every \(S\), let \(J(S)\) be the set of sieves of \(S\) covering for the
topology generated by the pretopology \(R\). Let \(J_R(S)\) be the subset
of \(J(S)\) consisting of the sieves defined by families in \(R(S)\).
Then \(J_R(S)\) is cofinal in \(J(S)\): every refinement of \(S\) contains
a sieve defined by a family in \(R(S)\).
\end{proposition}
\begin{proof}
For every \(S\), let \(J'(S)\) be the set of sieves of \(S\) containing a
sieve in \(J_R(S)\). One evidently has \(J'(S)\subset J(S)\).
To show that \(J(S)=J'(S)\), it suffices to show that the \(J'(S)\) form a
topology on \(\mathcal C\), that is, satisfy axioms (T 1) to (T 4).
Now (T 1), (T 3), (T 4) are evidently satisfied. It remains to verify (T 2).
\nde{In what follows, misprints in the original have been corrected.}

Thus let \(U\) be an element of \(J'(S)\) and \(C\) a sieve of \(S\);
one assumes that for every \(T\to U\), the sieve \(C\times_S T\) is in
\(J'(T)\), and one must prove that \(C\in J'(S)\).
By definition of \(J'\), \(U\) contains a refinement \(U'\) defined by a
family \(\{S_i\to S\}\in R(S)\). Since one has verified (T 3), it suffices
to prove that \(U'\cap C\in J'(S)\); one may therefore assume that \(U=U'\).
By hypothesis, for every \(i\), \(C\times_S S_i\in J'(S_i)\);
there therefore exists, for each \(i\), a covering family
\(\{T_{ij}\to S_i\}\in R(S_i)\) such that \(T_{ij}\to S_i\) factors through
\(C\times_S S_i\to S_i\). The morphism \(T_{ij}\to S\) therefore factors
through \(C\to S\), which shows that \(C\) contains the sieve defined by
the composite family \(\{T_{ij}\to S\}\), and one has finished by (P 2).
\end{proof}

Axioms (P 1) to (P 3) are those of [MA]. In view of the practical interest
of pretopologies, we shall interpret each important result with the aid
of a pretopology defining the given topology.

\begin{remark}{4.2.7}\label{IV.4.2.7}
One may introduce a slightly more general notion: for each \(S\) one gives
a set of covering families satisfying (P 1), (P 3) and Proposition 4.2.6.
This occurs in particular when the given families satisfy (P 1), (P 3)
and (C 5). The reader may consult [D].
\end{remark}

\begin{definition}{4.2.8}\label{IV.4.2.8}
Let \(\mathcal C\) be endowed with a topology, and let \(S\) be an object
of \(\mathcal C\). Let \(P(S')\) be a relation involving an argument
\(S'\in\Ob\mathcal C_{/S}\). One assumes that
\(\Hom(S'',S')\ne\varnothing\) implies \(P(S')\Rightarrow P(S'')\).
One says that \(P\) is \emph{true locally on \(S\)} for the topology considered
if the following equivalent conditions are satisfied:
\begin{enumeratei}
\item The set of \(S'\to S\) such that \(P(S')\) is true is a refinement of \(S\).
\item There exists a refinement of \(S\) such that \(P(S')\) is true for every
\(S'\) of this refinement.
\item (If the given topology is defined by a pretopology.) There exists a
covering family for this pretopology such that \(P(S')\) is true for every
\(S'\) of this family.
\end{enumeratei}
\end{definition}

\begin{example}{4.2.9}\label{IV.4.2.9}
Let \(f:X\to Y\) be an \(S\)-morphism. One will say that \(f\) is locally an
isomorphism if there exists a covering family \(\{S_i\to S\}\) such that,
for every \(i\), \(f\times_S S_i\) is an isomorphism.
It amounts to the same to require that there exist a refinement \(R\) of
\(S\) such that, for every \(T\to R\), \(X(T)\to Y(T)\) is an isomorphism.
\end{example}

One will see many other examples of \og local\fg{} language in what follows.

% typo: missing French preposition in "faisceau associé un préfaisceau".
\sgasection{4.3}{Presheaves, sheaves, sheaf associated to a presheaf}
\begin{definition}{4.3.1}\label{IV.4.3.1}
Let \(\mathcal C\) be a category. One calls a \emph{presheaf of sets} on
\(\mathcal C\) any contravariant functor from \(\mathcal C\) to the category
of sets. The category \(\widehat{\mathcal C}=\Hom(\mathcal C^\circ,\Ens)\)
is called the \emph{category of presheaves} on \(\mathcal C\).
If \(\mathcal C\) is endowed with a topology, one says that the presheaf
\(P\) is \emph{separated} (resp. is a \emph{sheaf}) if for every
\(S\in\Ob\mathcal C\) and every \(R\in J(S)\), the canonical map
\[
P(S)=\Hom(S,P)\longrightarrow\Hom(R,P)\tag{+}
\]
is injective (resp. bijective). One calls the full subcategory of
\(\widehat{\mathcal C}\) whose objects are the sheaves the
\emph{category of sheaves}, and denotes it by \(\widetilde{\mathcal C}\).
\nde{One will note that if \(Q\) is a subpresheaf of a separated presheaf
\(P\), then \(Q\) is separated. Indeed, for every sieve \(R\) of \(S\),
the composite map \(Q(S)\hookrightarrow P(S)\hookrightarrow P(R)\)
is injective and factors through \(Q(S)\to Q(R)\).}
\end{definition}

\begin{proposition}{4.3.2}\label{IV.4.3.2}
Let \(P\) be a separated presheaf (resp. a sheaf). For every functor
\(H\in\Ob\widehat{\mathcal C}\) and every \(R\in J(H)\), the canonical map
\[
\Hom(H,P)\longrightarrow\Hom(R,P)\tag{+}
\]
is injective (resp. bijective).
\end{proposition}
\begin{proof}
Indeed, let \(P\) be a separated presheaf, \(H\) a presheaf, \(R\in J(H)\),
and \(u,v:H\to P\) such that \(uj=vj\). For every \(f:S\to H\),
\(S\in\Ob\mathcal C\), \(R\times_H S\) is a refinement of \(S\) and
\(ufj_S=vfj_S\):
\[
\xymatrix@C=2.8em@R=2.3em{
 & R\ar[r]^j & H\ar[r]^{u,v} & P\\
 R\times_H S\ar[ur]^{f_R}\ar[r]_{j_S} & S\ar[ur]_f & & }.
\]
Since \(P\) is separated, one deduces \(uf=vf\). Since this is true for
every representable \(S\), one has \(u=v\).

Suppose now that \(P\) is a sheaf. Let \(g:R\to P\); let us show that it
factors through \(H\). For every \(f:S\to H\), \(S\in\Ob\mathcal C\),
\(g\circ f_R:R\times_H S\to P\) factors uniquely through \(S\), hence
defines a morphism \(h:S\to P\) which is evidently functorial in \(f\),
by uniqueness:
\[
\xymatrix@C=3.2em@R=2.5em{
R\times_H S\ar[r]^{f_R}\ar[d]_{j_S} & R\ar[r]^g\ar[d]^j & P\\
S\ar[r]_f\ar[urr]_h & H & }
\]
One has therefore defined, for every \(S\), a map from \(H(S)\) to \(P(S)\),
functorial in \(S\), hence a morphism from \(H\) to \(P\) which does satisfy
the required conditions.
\end{proof}

\begin{corollary}{4.3.2.1}\label{IV.4.3.2.1}
\nde{This corollary has been added.}
Let \(R,F\) be two sheaves. If \(R\) is a refinement of \(F\), then \(R=F\).
\end{corollary}
\begin{proof}
Indeed, suppose that \(R\) is a refinement of \(F\) and denote by \(j\)
the inclusion \(R\hookrightarrow F\). By 4.3.2, one has
\(\Hom(F,R)=\End(R)\), hence there exists \(\pi:F\to R\) such that
\(\pi\circ j=\mathrm{id}_R\). Similarly, one has
\(\End(F)=\Hom(R,F)\), and the equality \(j\circ\pi\circ j=j\)
implies \(j\circ\pi=\mathrm{id}_F\), so \(j\) is an isomorphism.
\end{proof}

\begin{proposition}{4.3.3}\label{IV.4.3.3}
\textup{([AS], 1.3).}
Let \(\mathcal C\) be a category. Let \(P\) be a presheaf on \(\mathcal C\);
for every \(S\in\Ob\mathcal C\), denote by \(J(S)\) the set of sieves \(R\)
of \(S\) such that for every \(T\to S\), the map
\[
\Hom(T,P)\longrightarrow\Hom(R\times_S T,P)\tag{+}
\]
is injective (resp. bijective). Then the \(J(S)\) define a topology on
\(\mathcal C\), i.e. satisfy axioms (T 1) to (T 4).
\end{proposition}

\begin{corollary}{4.3.4}\label{IV.4.3.4}
For every \(S\in\Ob\mathcal C\), let \(K(S)\) be a family of sieves satisfying
(T 1). Let \(P\) be a presheaf on \(\mathcal C\). For it to be separated
(resp. a sheaf) for the topology generated by the \(K(S)\), it is necessary
and sufficient that, for every \(S\in\Ob\mathcal C\) and every \(R\in K(S)\),
the canonical map
\[
\Hom(S,P)\longrightarrow\Hom(R,P)\tag{+}
\]
be injective (resp. bijective).
\end{corollary}

\begin{corollary}{4.3.5}\label{IV.4.3.5}
For each \(S\in\Ob\mathcal C\), let \(R(S)\) be a set of families of morphisms
of \(\mathcal C\) with target \(S\), satisfying (P 1) (for example defining a
pretopology). Let \(P\) be a presheaf on \(\mathcal C\). For \(P\) to be separated
(resp. a sheaf) for the topology generated by \(R\), it is necessary and sufficient
that for every \(S\in\Ob\mathcal C\) and every family \(\{S_i\to S\}\in R(S)\)
the map
\[
P(S)\longrightarrow\prod_i P(S_i)
\]
be injective (resp. the diagram
\[
P(S)\longrightarrow\prod_i P(S_i)
\rightrightarrows\prod_{i,j}P(S_i\times_S S_j)
\]
be exact).
\end{corollary}

\begin{definition}{4.3.6}\label{IV.4.3.6}
Let \(\mathcal C\) be a category. One calls the \emph{canonical topology}
on \(\mathcal C\) the finest topology for which all representable functors
are sheaves.
\end{definition}
\begin{corollary}{4.3.7}\label{IV.4.3.7}
For a sieve \(R\) of \(S\) to be a refinement for the canonical topology,
it is necessary and sufficient that, for every morphism \(T\to S\) of
\(\mathcal C\) and every \(X\in\Ob\mathcal C\), the canonical map
\[
\Hom(T,X)\longrightarrow\Hom(R\times_S T,X)
\]
be bijective.
\end{corollary}
\begin{definition}{4.3.8}\label{IV.4.3.8}
A covering sieve for the canonical topology will be called a
\emph{universal effective epimorphic sieve}.
\end{definition}
\begin{corollary}{4.3.9}\label{IV.4.3.9}
A universal effective epimorphic family defines a universal effective
epimorphic sieve. Conversely, every base-changeable family defining a
universal effective epimorphic sieve is universal effective epimorphic.
\end{corollary}

Let us return to the case where \(\mathcal C\) is endowed with an arbitrary
topology and pass to the construction of the sheaf associated to a presheaf
\(P\). Let \(S\) be an object of \(\mathcal C\). If \(R\supset R'\) are two
refinements of \(S\), one has a diagram
\[
\xymatrix{
\Hom(S,P)\ar[r]\ar[dr] & \Hom(R,P)\ar[d]\\
 & \Hom(R',P) }.
\]
The ordered set \(J(S)\) is filtered, as one has already remarked.
Since \(S\) is an element of \(J(S)\), one has an evident morphism
\[
\Hom(S,P)\longrightarrow\varinjlim_{R\in J(S)}\Hom(R,P).
\]
\begin{definition}{4.3.10.0}\label{IV.4.3.10.0}
% typo?: kept as printed: the N.D.E. has LP -> LQ for Q -> P.
\nde{The numbering 4.3.10.0 has been added for later references.
Moreover, it follows from the definition that if \(Q\to P\) is a
monomorphism, the same is true of \(LP\to LQ\); thus \(L\)
\og preserves monomorphisms\fg{} (see also 4.3.16 for a more general
result: \(L\) \og commutes with finite projective limits\fg).}
Set \(\check H^0(S,P)=\varinjlim_{R\in J(S)}\Hom(R,P)\).
One verifies that \(\check H^0(S,P)\) depends functorially on \(S\),
hence defines a functor \(LP\) by
\[
\Hom(S,LP)=\check H^0(S,P)=\varinjlim_{R\in J(S)}\Hom(R,P).\tag{++}
\]
\end{definition}
By construction one has morphisms
\[
\begin{gathered}
\ell_P:P\longrightarrow LP\\
z_R:\Hom(R,P)\longrightarrow\Hom(S,LP).
\end{gathered}
\]
\begin{lemma}{4.3.10}\label{IV.4.3.10}
(i) For every refinement \(R\) of \(S\) and every \(u:R\to P\),
the diagram
\[
\xymatrix{
P\ar[r]^{\ell_p} & LP\\
R\ar[u]^u\ar[r]_{i_R} & S\ar[u]_{z_R(u)} }
\]
is commutative.

(ii) For every morphism \(S\xrightarrow{v}LP\), there exists a refinement
\(R\) of \(S\) and a morphism \(u:R\to P\) with \(v=z_R(u)\).

(iii) Let \(Q\) be a functor and \(u,v:Q\to P\) such that
\(\ell_Pu=\ell_Pv\). Then the kernel of the pair \((u,v)\)
is a refinement of \(Q\).

(iv) Let \(u:R\to P\) and \(u':R'\to P\); for \(z_R(u)=z_{R'}(u')\),
it is necessary and sufficient that there exist a refinement
\(R''\subset R\cap R'\) of \(S\) such that \(u\) and \(u'\) agree on \(R''\).
\end{lemma}
\begin{proof}
(i): One must verify that \(z_R(u)i_R=\ell_Pu\). For this, it suffices to
verify that the composites of these two morphisms with every morphism
\(T\xrightarrow{g}R\), where \(T\) is representable, are equal.
% Proof continues in en-06.tex.
